% !TEX root = ../../main.tex
\section{Lưu trữ định danh nội dung (CAS) và Két sắt khóa cá nhân (BYOK)}
\label{sec:cas_byok}

\subsection{Nguyên lý lưu trữ định danh nội dung (CAS) và mã hóa xác thực AEAD}
\label{subsec:nguyen_ly_cas_byok}
Hạ tầng lưu trữ và an toàn dữ liệu của VeriScholar kết hợp hai kỹ thuật mã hóa hiện đại:

\begin{enumerate}[label=\textbf{(\arabic*)}]
    \item \textbf{Kho lưu trữ định danh nội dung (\textit{Content-Addressable Storage -- CAS}):}
    Trong hệ thống lưu trữ tệp thông thường, vị trí tệp được xác định bởi đường dẫn do người dùng đặt (dễ bị trùng lặp hoặc ghi đè). Ngược lại, CAS xác định vị trí tệp dựa trên chính nội dung nhị phân của nó thông qua hàm băm mật mã SHA-256:
    \begin{equation}
    \label{eq:sha256}
    H = \text{SHA-256}(\text{FileBytes})
    \end{equation}
    Tệp vật lý được lưu trữ phân tán theo cấu trúc cây thư mục 2 tầng dựa trên các ký tự đầu của chuỗi băm. Mỗi bản ghi tệp nhị phân được quản lý vòng đời bởi biến đếm tham chiếu (\textit{Reference Counting}). Khi một người dùng tải lên một tệp đã tồn tại trong hệ thống, hệ thống chỉ tăng biến đếm tham chiếu lên 1 mà không ghi thêm byte nào vào đĩa cứng.

    \item \textbf{Két sắt khóa cá nhân (\textit{BYOK}) và Chuẩn mã hóa AES-256-GCM:}
    Để bảo vệ khóa API cá nhân của người dùng (\textit{Bring Your Own Key -- BYOK}), hệ thống sử dụng thuật toán mã hóa đối xứng xác thực AEAD (\textit{Authenticated Encryption with Associated Data}) theo chế độ GCM (\textit{Galois/Counter Mode}) \cite{nguyendinhthuc2001}. 

    Thuật toán sử dụng khóa bí mật 256 bits ($K$) và vectơ khởi tạo ngẫu nhiên 96 bits ($\text{IV}$) cho mỗi lần mã hóa. Bản rõ khóa API được mã hóa bằng chế độ CTR (\textit{Counter}), đồng thời hàm băm vạn năng $\text{GHASH}$ kết hợp dữ liệu bổ sung không bảo mật ($\text{AAD}$) để tạo ra thẻ xác thực toàn vẹn 128 bits ($T$):
    \begin{equation}
    \label{eq:aes_gcm}
    C = \text{AES-CTR}_K(\text{Plaintext}), \quad T = \text{GHASH}_H(\text{AAD}, C) \oplus \text{AES}_K(\text{IV})
    \end{equation}
\end{enumerate}

\subsection[Luận giải lựa chọn CAS và AES-256-GCM cho VeriScholar]{Luận giải lựa chọn CAS và AES-256-GCM\\cho VeriScholar}
\label{subsec:luan_giai_cas_byok}
Đề tài khẳng định CAS và AES-256-GCM là hai nền tảng bảo mật và tối ưu hóa hạ tầng then chốt:
\begin{enumerate}[label=\textbf{(\arabic*)}]
    \item \textbf{Tiết kiệm 98\% chi phí lưu trữ tệp tài liệu học thuật:} Trong môi trường nghiên cứu đại học, nhiều sinh viên và giảng viên thường đọc chung các tài liệu kinh điển (ví dụ bài báo nền tảng về Transformer nặng gần 20MB \cite{vaswani2017}). Nhờ cơ chế khử trùng lặp CAS, hàng trăm người dùng có thể cùng nghiên cứu bài báo đó nhưng máy chủ chỉ lưu duy nhất một tệp vật lý trên đĩa cứng. Đồng thời, cấu trúc tệp tĩnh CAS hỗ trợ hoàn hảo giao thức HTTP Range Requests (\texttt{206 Partial Content}), cho phép trình duyệt tải luồng và hiển thị ngay lập tức từng trang tài liệu mà không cần chờ tải toàn bộ tệp PDF lớn.
    \item \textbf{Bảo mật cấp quân đội và chống tấn công thay đổi bản mã:} Chế độ AES-256-GCM bảo đảm cả tính bảo mật lẫn tính toàn vẹn. Nếu kẻ tấn công xâm nhập vào cơ sở dữ liệu và thay đổi dù chỉ 1 bit trong chuỗi mã hóa, hàm giải mã trên bộ nhớ RAM sẽ lập tức phát hiện thông qua thẻ xác thực $T$ và từ chối xử lý, ngăn chặn hoàn toàn các cuộc tấn công thay đổi bản mã (\textit{Bit-flipping}) hay tấn công đệm giải mã (\textit{Padding Oracle Attacks}).
    \item \textbf{Nguyên tắc không lưu vết trên RAM (\textit{RAM Zero-Retention}):} Khóa giải mã chỉ tồn tại tạm thời trong bộ nhớ RAM trong đúng thời gian thực thi lời gọi hàm AI và bị dọn dẹp ngay sau khi hoàn tất, đảm bảo quyền riêng tư tuyệt đối cho học giả.
\end{enumerate}
